%% SYNTHETIC TEST UNIT generated by writing/tools/make_outline_book.py from FINAL_OUTLINE_v2 -- not book text
\chapter{重夸克物理：$B$ 介子衰变与重夸克偶素}\label{ch:heavy}
本单元为测试用的合成文本（由 FINAL\_OUTLINE\_v2 自动生成），陈述内容从略。

\section{第 1 节（测试）}\label{sec:heavy-bdpi}

\begin{theorem}[$B\to D\pi$ 的全阶因子化 \tagA]\label{thm:heavy-bdpi}
（测试用合成陈述：内容从略。）
非微扰地位：\cref{asm:intro-bridge}。
\uses{def:excl-da, thm:scet-decoupling, asm:intro-bridge}
\end{theorem}
\begin{proofsketch}测试。\end{proofsketch}
\paragraph{地位与缺口}测试：本陈述的地位由\cref{rem:scet-assumptions}决定。

\begin{proposition}[重到轻形状因子：软重叠与硬旁观者 \tagA]\label{prop:heavy-btopi-ff}
（测试用合成陈述：内容从略。）
非微扰地位：\cref{asm:intro-bridge}。
\uses{def:excl-da, thm:scet-decoupling, asm:intro-bridge}
\end{proposition}
\begin{proofsketch}测试。\end{proofsketch}
\paragraph{地位与缺口}测试：本陈述的地位由\cref{prop:heavy-qcdf-endpoint}决定。

\begin{proposition}[无粲两体衰变的领头幂次 QCD 因子化 \tagP]\label{prop:heavy-qcdf-lp}
（测试用合成陈述：内容从略。）
非微扰地位：\cref{asm:intro-bridge}。
\uses{prop:heavy-btopi-ff, def:excl-da, asm:intro-bridge}
\end{proposition}
\paragraph{地位与缺口}测试：见文献。

\section{第 2 节（测试）}\label{sec:heavy-qcdf}

\begin{proposition}[幂次修正中的端点发散 \tagV]\label{prop:heavy-qcdf-endpoint}
（测试用合成陈述：内容从略。）
\uses{prop:heavy-qcdf-lp}
\end{proposition}
\paragraph{地位与缺口}测试：见文献。

\begin{theorem}[单举 $B$ 衰变的重夸克展开 \tagA]\label{thm:heavy-inclusive-b}
（测试用合成陈述：内容从略。）
非微扰地位：\cref{asm:intro-wightman}、\cref{asm:intro-composite}、\cref{asm:intro-af-coeff}、\cref{asm:intro-bridge}。
\uses{thm:dis-ope, thm:scet-decoupling, asm:intro-wightman, asm:intro-composite, asm:intro-af-coeff, asm:intro-bridge}
\end{theorem}
\begin{proofsketch}测试。\end{proofsketch}
\paragraph{地位与缺口}测试：本陈述的地位由\cref{thm:dis-ope}、\cref{rem:scet-assumptions}、\cref{conj:ee-duality}决定。

\begin{definition}[NRQCD 及其幂次计数]\label{def:heavy-nrqcd}
（测试用合成陈述：内容从略。）
\uses{def:prelim-qcd}
\end{definition}

\section{第 3 节（测试）}\label{sec:heavy-inclusive}

\begin{theorem}[衰变的 NRQCD 因子化 \tagA]\label{thm:heavy-nrqcd-decay}
（测试用合成陈述：内容从略。）
非微扰地位：\cref{asm:intro-bridge}、\cref{conj:ee-duality}。
\uses{def:heavy-nrqcd, asm:intro-bridge}
\end{theorem}
\begin{proofsketch}测试。\end{proofsketch}
\paragraph{地位与缺口}测试：见文献。

\begin{proposition}[匹配系数 \tagR]\label{prop:heavy-nrqcd-matching}
（$R_{\mathrm{PT}}$：在重整化微扰论中对红外正规化的部分子态严格。）
（测试用合成陈述：内容从略。）
\uses{def:heavy-nrqcd, def:prelim-msbar}
\end{proposition}

\begin{example}[阈值处的单圈 $c_v$]\label{ex:heavy-cv}
（测试用合成陈述：内容从略。）
\uses{prop:heavy-nrqcd-matching}
\end{example}

\section{第 4 节（测试）}\label{sec:heavy-nrqcd}

\begin{conjecture}[产生的 NRQCD 因子化 \tagC]\label{conj:heavy-nrqcd-production}
（测试用合成陈述：内容从略。）
\uses{def:heavy-nrqcd}
\end{conjecture}
\paragraph{地位与缺口}测试：见文献。

\begin{proposition}[大 $p_T$ 重夸克偶素产生 \tagA]\label{prop:heavy-quarkonium-fragmentation}
（测试用合成陈述：内容从略。）
非微扰地位：\cref{asm:intro-bridge}。
\uses{def:pdf-ff, asm:intro-bridge}
\end{proposition}
\begin{proofsketch}测试。\end{proofsketch}
\paragraph{地位与缺口}测试：本陈述的地位由\cref{conj:dy-coloured-final}、\cref{conj:heavy-nrqcd-production}决定。
